\reading{LTR-4}{The Investment Function in Financial-Services Management}{DEFER}{1}
\markboth{LTR-4\ \ The Investment Function}{}

\begin{lrkeybox}[Learning objectives — official 2026 spine wording]
\begin{enumerate}[label=\textbf{\alph*.}]
\item Compare various money market and capital market instruments and discuss their advantages and disadvantages.
\item Identify and discuss various factors that affect the choice of investment securities by a bank.
\item Apply investment maturity strategies and maturity management tools based on the yield curve and duration.
\end{enumerate}
{\footnotesize\color{FRMGrey}Source: Peter Rose and Sylvia Hudgins, \textit{Bank Management \& Financial Services}, 9th Edition, Chapter 10.}
\end{lrkeybox}

% =====================================================================
\lo{LTR-4 a}{Compare various money market and capital market instruments and discuss their advantages and disadvantages.}

\begin{lrdefbox}[The classification, and the line between them]
\term{Money market instruments} reach maturity \textbf{within one year} and are noted for their low risk and ready marketability. \term{Capital market instruments} have remaining maturities \textbf{beyond one year} and are generally noted for their higher expected rate of return and capital gains potential.
\end{lrdefbox}

\subsection{Money market instruments}

\begin{center}
\scriptsize
\begin{longtable}{@{}p{2.9cm} p{5.3cm} p{5.3cm}@{}}
\toprule
\rowcolor{FRMSoft}\textbf{Instrument} & \textbf{Key advantages} & \textbf{Key disadvantages} \\
\midrule
\endfirsthead
\toprule
\rowcolor{FRMSoft}\textbf{Instrument} & \textbf{Key advantages} & \textbf{Key disadvantages} \\
\midrule
\endhead
\textbf{Treasury bills} \newline {\color{FRMGrey}notes rate $\star\star\star$} &
Safety and high liquidity; good collateral for borrowing; can be pledged behind government deposits. &
Low yields relative to other financial instruments; taxable income. \\
\addlinespace
\textbf{Short-term Treasury notes and bonds} \newline {\color{FRMGrey}$\star\star\star$} &
Safety, good resale market; good collateral for borrowing; offer yields usually higher than T-bill yields. &
More price risk than T-bills; taxable gains and income. \\
\addlinespace
\textbf{Federal agency securities} \newline {\color{FRMGrey}$\star\star$} &
Safety, good to average resale market; good collateral for borrowing; higher yields than on US government securities. &
Less marketable than Treasury securities; taxable gains and income. \\
\addlinespace
\textbf{Certificates of deposit} \newline {\color{FRMGrey}$\star\star$} &
Insured to at least \$100,000; yields higher than on T-bills; large denominations often marketable through dealers. &
Limited resale market on longer-term CDs; taxable income. \\
\addlinespace
\textbf{International Eurocurrency deposits} \newline {\color{FRMGrey}$\star\star$} &
Low risk; higher yields than on many domestic CDs. &
Volatile interest rates; taxable income. \\
\addlinespace
\textbf{Bankers' acceptances} \newline {\color{FRMGrey}$\star$} &
Low risk due to multiple credit guarantees. &
Limited availability at specific maturities; issued in odd denominations; taxable income. \\
\addlinespace
\textbf{Commercial paper} \newline {\color{FRMGrey}$\star$} &
Low risk due to high quality of borrowers. &
Volatile market; poor resale market; taxable income. \\
\addlinespace
\textbf{Short-term municipal obligations} \newline {\color{FRMGrey}$\star$} &
Tax-exempt interest income. &
Limited resale market; taxable capital gains. \\
\bottomrule
\end{longtable}
\end{center}

\subsection{Capital market instruments}

\begin{center}
\scriptsize
\begin{longtable}{@{}p{2.9cm} p{5.3cm} p{5.3cm}@{}}
\toprule
\rowcolor{FRMSoft}\textbf{Instrument} & \textbf{Key advantages} & \textbf{Key disadvantages} \\
\midrule
\endfirsthead
\toprule
\rowcolor{FRMSoft}\textbf{Instrument} & \textbf{Key advantages} & \textbf{Key disadvantages} \\
\midrule
\endhead
\textbf{Treasury notes and bonds} \newline {\color{FRMGrey}$\star\star$} &
Safety, good resale market; good collateral for borrowing; may be pledged behind government deposits. &
Low yields relative to long-term private securities; taxable gains and income; limited supply of longest-term issues. \\
\addlinespace
\textbf{Municipal (state and local government) bonds} \newline {\color{FRMGrey}$\star\star$} &
Tax-exempt interest income; high credit quality; liquidity and marketability of selected securities. &
Volatile market, some issues have limited resale potential; taxable capital gains. \\
\addlinespace
\textbf{Corporate notes and bonds} \newline {\color{FRMGrey}$\star$} &
Higher pretax yields than on government securities; aid in locking in higher long-term rates of return. &
Limited resale market; inflexible terms; taxable gains and income. \\
\addlinespace
\textbf{Asset-backed securities} \newline {\color{FRMGrey}$\star$} &
Higher pretax yields than on Treasury securities; collateral for borrowing additional funds. &
Less marketable and more unstable in price than Treasury securities; may carry substantial default risk; taxable gains and income. \\
\bottomrule
\end{longtable}
\end{center}
\figcap{Star ratings are the source notes' own importance grading, retained. Read the disadvantage column for the tax word: \emph{taxable income} appears on almost every line, and the two municipal entries are the exceptions — tax-exempt \textbf{interest}, but taxable \textbf{capital gains}. That split is the discriminator.}

\begin{lrtrapbox}[Two discriminators that carry the whole objective]
\begin{itemize}
\item \textbf{Municipals are half-exempt.} Interest income is tax-exempt; capital gains are taxable. A distractor claiming municipals are fully tax-exempt is the obvious build.
\item \textbf{Federal agency securities yield \emph{more} than US government securities and are \emph{less} marketable than Treasuries.} Both halves are directional and both get flipped.
\end{itemize}
\end{lrtrapbox}

% =====================================================================
\lo{LTR-4 b}{Identify and discuss various factors that affect the choice of investment securities by a bank.}

\subsection{Investment security selection and risks}

\begin{enumerate}
\item \textbf{Expected rate of return.} Banks aim for securities with high after-tax yields, considering factors such as YTM, HPY, and tax implications.
  \begin{enumerate}[label=\alph*)]
  \item \term{Yield to maturity (YTM)}: best for most securities, except equities.
  \item \term{Holding period yield (HPY)}: for securities sold before maturity.
  \end{enumerate}
\item \textbf{Tax exposure.} Banks favour tax-exempt municipal bonds, and securities with a \emph{lower tax burden}, compared to taxable options.
  \begin{enumerate}[label=\alph*)]
  \item \term{Tax swapping}: selling low-yielding securities at a loss to offset income, and buying a higher-yielding security.
  \item \term{Portfolio shifting}: selling securities at a loss to reduce tax liability.
  \end{enumerate}
\end{enumerate}

\begin{lrexbox}[After-tax return on a bank-qualified municipal bond]
A bank purchases a municipal bond with a 6\% gross nominal rate of return. The bank's borrowing cost to purchase the bond is 5\%, and the bank is in the top income tax bracket of 35\%. The bond is a bank-qualified bond, and 80\% of the interest expense is tax deductible. Calculate the net after-tax return on the bond.
\tcblower
\textbf{Step 1 — the raw spread.} The bond earns 6\% and the funding costs 5\%:
\[ 6\% - 5\% = 1\% \]
This 1\% is the pre-tax carry. It is \emph{not} the answer, and it is a natural distractor.

\textbf{Step 2 — value the interest deduction.} Only 80\% of the 5\% interest expense is deductible, and the deduction is worth the 35\% tax rate:
\[ 0.35 \times 0.80 \times 5\% = 1.4\% \]

\textbf{Step 3 — add the shield to the carry.}
\[ \text{Net after-tax return} = 6\% - 5\% + (0.35 \times 0.80 \times 5\%) = \mathbf{2.4\%} \]

\textbf{What this means.} The tax shield on the \emph{funding} is worth more than the carry itself. The municipal's interest is already tax-exempt, so no tax is taken off the 6\%; the only tax term in the calculation runs the other way, as a deduction against the cost of carrying it. That asymmetry is the entire reason bank-qualified municipals exist.
\end{lrexbox}

\begin{lrtrapbox}[The sign on the tax term]
The tax term is \textbf{added}, not subtracted. It is a deduction on an expense, not a tax on income. A distractor computing $6\% - 5\% - 1.4\% = -0.4\%$ is the sign-flip build, and $1\%$ is the dropped-term build.
\end{lrtrapbox}

\subsection{Other risks affecting the choice}

\begin{center}
\small
\begin{tabularx}{\textwidth}{@{}p{3.5cm} X@{}}
\toprule
\rowcolor{FRMSoft}\textbf{Risk} & \textbf{What it is} \\
\midrule
\textbf{Interest rate risk} & Impacts portfolio value and loan demand. \\
\textbf{Credit risk} & Risk of the issuer defaulting on payments. Credit ratings measure relative credit risk. \\
\textbf{Business risk} & Risk of a weakening local economy. \\
\textbf{Liquidity risk} & Risk of difficulty selling a security quickly. \\
\textbf{Call risk} & The issuer's right to repurchase bonds at a predetermined price — unfavourable to the investor when interest rates \textbf{decline}. \\
\textbf{Prepayment risk} & Risk of not receiving expected cash flows, often in mortgage-backed securities (mortgage refinancing). \\
\textbf{Inflation risk} & Erosion of investment value due to rising prices. \\
\textbf{Pledging requirements} & Collateral needed for government deposits. \\
\bottomrule
\end{tabularx}
\end{center}

\begin{lrtrapbox}[Call risk bites when rates fall]
The issuer calls when it can refinance more cheaply, which is when rates have \textbf{declined}. Call risk and prepayment risk are the same mechanism on different instruments, and both are falling-rate events. A distractor pairing call risk with rising rates is the polarity build.
\end{lrtrapbox}

% =====================================================================
\lo{LTR-4 c}{Apply investment maturity strategies and maturity management tools based on the yield curve and duration.}

The maturity of investments held is an important consideration for financial institutions. Institutions look at whether short- or long-term maturities are more appropriate, and what the distribution of maturities should be.

\subsection{The five investment maturity strategies}

\begin{center}
\scriptsize
\begin{longtable}{@{}p{2.9cm} p{5.0cm} p{5.6cm}@{}}
\toprule
\rowcolor{FRMSoft}\textbf{Strategy} & \textbf{What the portfolio looks like} & \textbf{Advantage} \\
\midrule
\endfirsthead
\toprule
\rowcolor{FRMSoft}\textbf{Strategy} & \textbf{What the portfolio looks like} & \textbf{Advantage} \\
\midrule
\endhead
\textbf{The Ladder, or Spaced-Maturity Policy} \newline {\color{FRMGrey}$\star\star\star$} &
Divide the investment portfolio \textbf{equally among all maturities} acceptable to the investing institution — for example 20\% of the portfolio in each of years 1 to 5. &
Reduces investment income fluctuations; requires little management expertise. \\
\addlinespace
\textbf{The Front-End Load Maturity Policy} \newline {\color{FRMGrey}$\star\star\star$} &
All security investments are \textbf{short-term} — for example 30\% at 1 year and 70\% at 2 years, with nothing beyond. &
Strengthens the financial firm's liquidity position and avoids large capital losses if market interest rates \textbf{rise}. \\
\addlinespace
\textbf{The Back-End Load Maturity Policy} \newline {\color{FRMGrey}$\star\star\star$} &
All security investments are \textbf{long-term} — for example everything spread across years 5 to 10. &
Maximizes income potential from security investments if market interest rates \textbf{fall}. \\
\addlinespace
\textbf{The Barbell Strategy} \newline {\color{FRMGrey}$\star\star\star$} &
Security holdings are \textbf{divided between short-term and long-term}, with half or a substantial percentage at each end and nothing in the middle. &
Helps meet liquidity needs with the short-term securities and achieve earnings goals from the higher potential earnings of the long-term portion. It also gains a large \textbf{convexity} benefit. \\
\addlinespace
\textbf{The Rate Expectations Approach} \newline {\color{FRMGrey}$\star\star\star$} &
\textbf{Change the mix} of investment maturities as the interest-rate outlook changes: shift toward long-term securities if rates are expected to \textbf{fall}, toward short-term securities if rates are expected to \textbf{rise}. &
Maximizes the potential for earnings --- and also for losses. \\
\bottomrule
\end{longtable}
\end{center}

\begin{lrkeybox}[The decision grid the source notes give for choosing between them]
\begin{center}
\small
\begin{tabularx}{0.95\textwidth}{@{}p{5.6cm} X@{}}
\toprule
\rowcolor{FRMSoft}\textbf{What the institution has} & \textbf{Strategy} \\
\midrule
No opinion, no money & The Ladder, or Spaced-Maturity Policy \\
Liquidity matters & The Front-End Load Maturity Policy \\
Return matters & The Back-End Load Maturity Policy \\
Both liquidity and return matter & The Barbell Strategy \\
Money and ideas & The Rate Expectations Approach \\
\bottomrule
\end{tabularx}
\end{center}
This mapping is the fastest route into a scenario question. The stem tells you which of liquidity, return, opinion and expertise the institution has; the grid names the strategy.
\end{lrkeybox}

\begin{lrtrapbox}[Front-end versus back-end: get the rate direction right]
\begin{itemize}
\item \textbf{Front-end} is all short. Short maturities mean small price losses when rates \textbf{rise}. Its advantage is stated against a \emph{rising} rate scenario.
\item \textbf{Back-end} is all long. Long maturities mean large price gains when rates \textbf{fall}. Its advantage is stated against a \emph{falling} rate scenario.
\item \textbf{Rate expectations} shifts \emph{toward long} when rates are expected to fall and \emph{toward short} when rates are expected to rise — the same logic, applied dynamically.
\end{itemize}
All three statements are polarity-sensitive and all three are offered reversed. Say the direction and the maturity out loud together.
\end{lrtrapbox}

\subsection{Maturity management tools}

The two primary maturity management tools are the \textbf{yield curve} and \textbf{duration}.

\begin{lrdefbox}[The yield curve]
The yield curve is simply a picture of how market interest rates differ across loans and securities of varying term to maturity.
\end{lrdefbox}

\begin{lrkeybox}[Riding the yield curve]
\begin{itemize}
\item The strategy is based on the fact that as time passes, the bond's remaining maturity and duration decrease.
\item When the \textbf{yield curve} is upward sloping, buy long-term bonds and sell short-term bonds.
\item Benefit from: higher gain during price appreciation and lower loss during price depreciation. Then, after the yield declines, the manager sells the bond and rolls out the curve to repeat the process, buying another bond at the end of a steep segment of the curve.
\end{itemize}
\end{lrkeybox}

\subsection{Pursuing the carry trade}

A carry trade involves buying a security and financing it at a rate that is lower than the security's yield.

\begin{lrdefbox}[Duration]
Duration is an essential tool in investment analysis, measuring an investment's sensitivity to interest-rate changes.
\end{lrdefbox}

\begin{lrtrapbox}[Consolidated trap summary — LTR-4]
\begin{itemize}
\item \textbf{Sibling.} The \textbf{Ladder} spreads the portfolio evenly across all maturities. The \textbf{Back-End Load} is the all-long strategy. The source notes' Crash layer swaps these two, and the swap is corrected here — see the corrections record.
\item \textbf{Polarity.} Front-end load protects against \emph{rising} rates. Back-end load profits from \emph{falling} rates. Rate expectations shifts long when rates are expected to fall.
\item \textbf{Polarity.} Call risk and prepayment risk are \emph{falling}-rate events.
\item \textbf{Scope.} Municipal interest is tax-exempt; municipal \emph{capital gains} are taxable.
\item \textbf{Sign.} In the after-tax return calculation the tax term is \textbf{added} (a deduction on the funding cost), not subtracted.
\item \textbf{Intermediate result.} The 1\% raw carry and the 1.4\% tax shield are both natural offered distractors against the correct 2.4\%.
\item \textbf{Definition.} Riding the yield curve requires an \emph{upward-sloping yield curve}, not an upward-sloping price.
\item \textbf{Sibling.} YTM applies to most securities except equities; HPY applies to securities \emph{sold before maturity}.
\end{itemize}
\end{lrtrapbox}
